% TOP_PROBLEM: {"id": 30007199, "title": "Recognizing virtually polycyclic groups from quasi-isometry", "category_id": 17, "category": {"id": 17, "name": "group_theory", "display_name": "Group Theory", "description": "Problems about groups, group actions, representations, and related algebraic structures.", "slug": "group-theory", "order_index": 17, "created_at": "2026-07-31T15:26:25.671Z"}, "status": "open"}
% FIRST_POSTED: 2026-10-06
% ATTEMPT_STATUS: solved
% OPENAI_CLAIM: {"schema_version":1,"source_repo":"https://github.com/openai/math","source_commit":"adc7f1241b42e322a6451854ab7e4b4c146bf78a","release_date":"2026-10-06","match":"full","resolution":"proved","independently_reviewed":false,"second_pass":"agreed","families":[{"family":"255","title":"Quasi-isometric recognition of virtually polycyclic groups","subject":"Group theory","manuscripts":[{"title":"Quasi-isometric Recognition of Virtually Polycyclic Groups","pdf_path":"preprints/quasi-isometric-recognition-of-virtually-polycyclic-groups-September-24-2026/paper.pdf","readme_path":"preprints/quasi-isometric-recognition-of-virtually-polycyclic-groups-September-24-2026/README.md","date":"2026-09-24","theorem_ref":"thm:recognition (introduction.tex:10)","statement_sources":[{"path":"preprints/quasi-isometric-recognition-of-virtually-polycyclic-groups-September-24-2026/build/introduction.tex","line":10,"label":"thm:recognition"}]}],"lean":{"doc_path":"lean/docs/255.md","comparators":["lean/ComparatorChallenges/PolycyclicRecognition.json","lean/ComparatorChallenges/PolycyclicRecognition.lean"],"declarations":[],"covers_main_theorem":true},"reasoning_trace":null}]}
\documentclass[11pt]{article}
\usepackage[margin=25mm]{geometry}
\usepackage{fontspec}
\setmainfont{Latin Modern Roman}
\usepackage{amsmath,amssymb,mathtools}
\usepackage{unicode-math}
\setmathfont{Latin Modern Math}
\usepackage{enumitem}
\usepackage{xurl,hyperref,bookmark}
\hypersetup{colorlinks=true,urlcolor=blue}
\setcounter{secnumdepth}{0}
\setlength{\parindent}{0pt}
\setlength{\parskip}{5pt}
\setlength{\emergencystretch}{3em}
\begin{document}
\section*{Recognizing virtually polycyclic groups from quasi-isometry}

\subsection{OpenAI's claim}
OpenAI presents a proof in the selected manuscript. Virtual polycyclicity is claimed determined by the coarse geometry among groups with finite generators.

\subsection{Mathematical statement of the claimed result}
\paragraph{Result 1.} Let \(P\) be finitely generated and virtually polycyclic, and let \(H\) be any finitely generated group. The source claims
\[
 H\text{ quasi-isometric to }P\quad\Longrightarrow\quad H\text{ virtually polycyclic}.
\]
Quasi-isometry refers to word metrics from finite generating sets.

\subsection{Outline of the argument}
Cartan models and weighted boxes define canonical coarse height data. Path packing and recurrent switches force stationary slopes, and neighboring affine chart limits agree; the resulting solvable-lattice structure identifies virtual polycyclicity without fixing the original ambient Lie group.

\subsection{Scope relative to this problem}
The recognition theorem quantifies every finitely generated H and finitely generated virtually polycyclic P and has exactly the requested quasi-isometry implication. lean/docs/255.md documents the exact recognition statement. Within these assumptions, this is a claimed affirmative answer to the stated question.

\subsection{Formal verification}
For family 255, the documented theorem recognizes virtual polycyclicity by quasi-isometry and places a finite-index subgroup in a possibly different solvable Lie lattice; a height-control supplement is also listed. The listed Comparator files are \texttt{PolycyclicRecognition.lean}. The inventory does not supply a named declaration for these Comparator entries. This site has not executed the Lean code. The scope descriptions do not certify the remaining manuscript arguments.

\subsection{Status caveat}
OpenAI attributes these manuscripts to work by an internal model. Its release cautions that arguments without formalization can contain mistakes. This catalogue has not independently checked the proofs; the labels report the mathematical claims attributed to OpenAI.

\subsection{Sources}
\begin{itemize}
\item OpenAI, \emph{Quasi-isometric Recognition of Virtually Polycyclic Groups}, 2026-09-24; citation key \texttt{OAI:quasi-isometric-recognition-of-virtually-polycyclic-groups-September-24-2026}. \href{https://github.com/openai/math/blob/adc7f1241b42e322a6451854ab7e4b4c146bf78a/preprints/quasi-isometric-recognition-of-virtually-polycyclic-groups-September-24-2026/paper.pdf}{PDF (GitHub)}; \href{https://github.com/openai/math/raw/adc7f1241b42e322a6451854ab7e4b4c146bf78a/preprints/quasi-isometric-recognition-of-virtually-polycyclic-groups-September-24-2026/paper.pdf}{PDF (direct)}; \href{https://github.com/openai/math/blob/adc7f1241b42e322a6451854ab7e4b4c146bf78a/preprints/quasi-isometric-recognition-of-virtually-polycyclic-groups-September-24-2026/README.md}{Manuscript page and citation}; \href{https://github.com/openai/math/blob/main/preprints/quasi-isometric-recognition-of-virtually-polycyclic-groups-September-24-2026/paper.pdf}{Latest version}; Statement source: \href{https://github.com/openai/math/blob/adc7f1241b42e322a6451854ab7e4b4c146bf78a/preprints/quasi-isometric-recognition-of-virtually-polycyclic-groups-September-24-2026/build/introduction.tex\#L10}{thm:recognition}.
\item Family 255: \href{https://github.com/openai/math/blob/adc7f1241b42e322a6451854ab7e4b4c146bf78a/lean/docs/255.md}{Lean scope}.
\item \href{https://github.com/openai/math/blob/adc7f1241b42e322a6451854ab7e4b4c146bf78a/lean/ComparatorChallenges/PolycyclicRecognition.json}{Comparator PolycyclicRecognition.json}.
\item \href{https://github.com/openai/math/blob/adc7f1241b42e322a6451854ab7e4b4c146bf78a/lean/ComparatorChallenges/PolycyclicRecognition.lean}{Comparator PolycyclicRecognition.lean}.
\item \href{https://github.com/openai/math/blob/adc7f1241b42e322a6451854ab7e4b4c146bf78a/CONTENTS.md}{OpenAI catalogue entry}.
\end{itemize}
\end{document}
